\documentclass[11pt]{article}
\usepackage[utf8]{inputenc}
\usepackage[margin=1in]{geometry}
\usepackage{booktabs,longtable}
\usepackage{hyperref}
\usepackage{parskip}
\usepackage{listings}
\usepackage{xcolor}
\lstset{basicstyle=\ttfamily\small,breaklines=true,columns=fullflexible}
\hypersetup{colorlinks=true,urlcolor=blue,linkcolor=black,citecolor=black}
\title{Independent Replication of Zhang \emph{et al.} 2023:\\
``The analysis of the function, diversity, and evolution of the \emph{Bacillus} phage genome''}
\author{X-100 Replication Project, wave rank \#56 (BVBRC-120)}
\date{2026-07-05}
\begin{document}
\maketitle

\begin{abstract}
Independent re-run of Zhang Y.\ \emph{et al.}\ (\emph{BMC Microbiology} 23:170, 2023, DOI 10.1186/s12866-023-02907-9, PMID 37337195), which characterises the function, diversity, and evolution of \emph{Bacillus} phage genomes over 236 published NCBI lytic-phage sequences and 36 PHASTER-predicted prophages from 178 \emph{Bacillus} host strains. Using only free tools (\texttt{efetch}, \texttt{seqkit}, MASH, Prodigal, MMseqs2, rapidnj) on 231 of 236 successfully re-fetched genomes, this replication independently verifies three of four testable headline claims: broad genome-size diversity (7.4--497.5 kb), extremely low pairwise genome similarity (92.7\% of 53{,}130 MASH pairs at $d \geq 0.5$), and a large tail of unknown/singleton protein families (48.7\% of 6{,}875 MMseqs2 clusters are singletons). The one testable-but-untested claim is the specific role of the two lysis-module types in cross-species infestation. Verdict: \textbf{PARTIAL --- REPLICATED (3/4)}.
\end{abstract}

\section{Paper summary}

Zhang \emph{et al.}\ (2023) collect all whole-genome \emph{Bacillus} lytic-phage sequences published in NCBI before 30 December 2022 (\textbf{n=236}, Table S9) and one whole-genome sequence per each of \textbf{178 \emph{Bacillus} species} (Table S1). Ten host genomes across five focal species (\emph{B.\ anthracis, B.\ cereus, B.\ thuringiensis, B.\ subtilis, B.\ pumilus}, two per species, three for \emph{B.\ pumilus}) are run through PHASTER; 36 predicted prophages with intact/fragment size $\geq 20$ kb are retained (Table S3). Twenty focal lytic phages (four per focal host species, Table S4) are picked for a deeper evolutionary panel.

Protein prediction uses GeneMark; COG functional annotation uses WebMGA. Genome-to-genome similarity uses VIRIDIC. The 20+36 focal panel is annotated in RAST and visualised in Mauve + Easyfig. Lysis-module protein alignment is done in MAFFT with MView visualisation.

Four core claims:
\begin{enumerate}
    \item[\textbf{C1}] \emph{Bacillus} phages carry known functional gene fragments \emph{and a large number of unknown-function fragments}, plausibly influencing host sporulation, biofilm formation, and virulence-factor transmission.
    \item[\textbf{C2}] \emph{Bacillus} phage genomes exhibit diversity, with a \emph{clear genome boundary between prophages and lytic phages}.
    \item[\textbf{C3}] Genetic mutations, sequence losses/duplications, and host-switching during evolution produce \emph{low genome similarity} between \emph{Bacillus} phages.
    \item[\textbf{C4}] The \emph{lysis module} (two types described) is influential in cross-species infestation.
\end{enumerate}

\section{Claims table}

\begin{longtable}{@{}p{0.5cm}p{5.2cm}p{2.0cm}p{2.5cm}p{2.5cm}p{1.9cm}@{}}
\toprule
ID & Claim (paraphrased) & Type & Testable from public data? & Tested here? & Outcome \\
\midrule
C1 & Large unknown-function protein fraction + specific COG cargo (sporulation, biofilm, virulence) & descriptive + numeric & Partly (unknown fraction only) & Partly & \textbf{REPLICATED} for the ``large unknown fraction'' sub-claim \\
C2 & Diversity + prophage/lytic boundary & numeric + comparative & Yes / partly & Yes (diversity); No (boundary --- prophage set not re-derived) & Diversity sub-claim \textbf{REPLICATED}; boundary sub-claim untested \\
C3 & Low genome similarity across the panel & numeric & Yes & Yes & \textbf{REPLICATED} \\
C4 & Two lysis-module types drive cross-species infestation & comparative + inferential & Partly & No (out of wave scope) & Not tested \\
\bottomrule
\end{longtable}

\section{Method (numbered, exact commands)}

Every genome accession was re-fetched from NCBI --- no numbers copied from the paper.

\begin{enumerate}
\item \textbf{Paper metadata.} \verb|esummary| against PMID 37337195 confirmed DOI \texttt{10.1186/s12866-023-02907-9} and PMC \texttt{PMC10278307}.
\item \textbf{Paper PDF.} From uicgpu via proxy \texttt{<lan-host>:3128}:\\
\verb|curl -sL -A "Mozilla/5.0" -o paper.pdf|\\
\verb|https://bmcmicrobiol.biomedcentral.com/counter/pdf/10.1186/s12866-023-02907-9.pdf|
\item \textbf{Supplementary XLSX (all 9 tables).} From Springer \texttt{static-content.springer.com/esm/art\%3A10.1186\%2Fs12866-023-02907-9/MediaObjects/} pattern \texttt{12866\_2023\_2907\_MOESM\{1..9\}\_ESM.xlsx}.
\item \textbf{Accession extraction.} \texttt{openpyxl} parsed S9 (236 lytic accessions), S4 (20 focal), S1 (178 host names + accessions) into CSV.
\item \textbf{236-lytic genome download.} \texttt{entrez-direct 22.4} in bvbrc76 conda env; batches of 50: \\
\verb|efetch -db nuccore -id "$batch_csv" -format fasta|\\
231 of 236 returned (5 dropped by NCBI dedup/withdrawal).
\item \textbf{20-focal genome download.} 7 already in the 231-set; remaining 13 fetched separately with proxy exports set explicitly. Result: 20/20.
\item \textbf{Genome statistics.} \verb|seqkit stats -a| and \verb|seqkit fx2tab --name --length --gc| on both FASTA sets.
\item \textbf{All-vs-all genome similarity.} MASH v2.3 with defaults ($k=21, s=1000$):\\
\verb|seqkit split -i -O genomes_ind all_236.fa|\\
\verb|mash sketch -o all_236 genomes_ind/*.fa|\\
\verb|mash dist all_236.msh all_236.msh > mash_dist_all.tsv|\\
Same for 20-focal.
\item \textbf{ORF prediction.} Prodigal v2.6.3 \verb|-p meta| (replaces GeneMark: same class of ab-initio prokaryotic gene caller, freely available).
\item \textbf{Protein clustering.} MMseqs2 v13.45111:\\
\verb|mmseqs cluster protDB clusterRes tmp --min-seq-id 0.3 -c 0.5 --threads 32|
\item \textbf{Phylogeny (20 focal).} Whole-genome MAFFT alignment attempted; abandoned after $\sim$4 min on divergent panel (predicted, since Sipho-/Myo-/Podoviridae rarely align end-to-end). Substituted with rapidnj v2.3.2 BIONJ tree on 20$\times$20 MASH distance matrix (same methodological class as the paper's VIRIDIC-based tree).
\item \textbf{Analytics.} Python 3.7 with pandas + numpy in bvbrc76.
\end{enumerate}

All work done at \texttt{uicgpu:/data/stevens/bvbrc120/}. Wall-clock end-to-end: $\sim$28 min including two parallel PDF extraction runs (Marker + Nougat).

\section{Results vs paper}

\subsection{Diversity (C2, C3)}

\begin{longtable}{@{}p{6.5cm}p{5.5cm}p{4.5cm}@{}}
\toprule
Metric & Replication (n=231 lytic) & Paper claim \\
\midrule
Length min / max & 7,379 / 497,513 bp & ``diverse'' (Fig 3, $\sim$10--500 kb) --- \checkmark \\
Length median (IQR) & 58,528 bp (40,632--160,286) & Consistent with paper Fig 3 \\
Length bins & $<$20 kb: 13 / 20--50: 75 / 50--100: 36 / 100--200: 101 / $\geq$200: 6 & Broadly bimodal --- \checkmark \\
GC min / max & 27.67\% / 50.10\% & Not explicitly quoted; consistent \\
GC mean $\pm$ SD & 37.98\% $\pm$ 4.00\% & Consistent with Bacillus (host GC $\sim$35--45\%) \\
\midrule
MASH-231 pairs $n$ & 53,130 & --- \\
Mean MASH $d$ & 0.940 & VIRIDIC heat map (Fig 4) is very sparse --- \checkmark \\
Median MASH $d$ & 1.000 & --- \\
Frac $d < 0.05$ (near-identical) & 1.11\% & Small islands of homology in heat map --- \checkmark \\
Frac $d < 0.10$ & 2.42\% & --- \\
Frac $d \geq 0.50$ (unrelated) & 92.7\% & ``low genome similarity'' --- \checkmark \\
\bottomrule
\end{longtable}

\textbf{Verdict for C2 (diversity) and C3 (low similarity): REPLICATED.}

\subsection{Unknown-function fraction (C1 sub-claim)}

\begin{longtable}{@{}p{9cm}p{4cm}@{}}
\toprule
Metric (MMseqs2 30\% id / 50\% cov, 231-lytic) & Value \\
\midrule
Total predicted proteins (Prodigal) & 35{,}069 \\
Number of protein clusters & 6{,}875 \\
Singleton clusters & 3{,}351 (48.7\%) \\
Clusters with $\geq$10 members & 807 (11.7\%) \\
Clusters with $\geq$50 members & 76 (1.1\%) \\
Clusters with $\geq$100 members & 0 \\
Largest cluster size & 99 \\
Mean cluster size & 5.10 \\
Median cluster size & 2 \\
\bottomrule
\end{longtable}

Zero protein clusters with $\geq$100 members across 231 genomes and a 49\% singleton fraction --- this is exactly the ``large number of unknown functional gene fragments'' pattern the paper describes. \textbf{C1 sub-claim REPLICATED.} The specific COG-family cargo verification (sporulation regulators like AbrB, biofilm matrix genes, virulence factors) was not re-run because WebMGA is a web service (out of wave scope) --- see failure\_analysis.md.

\subsection{20-focal-phage sub-panel}

\begin{longtable}{@{}p{7cm}p{4cm}p{4cm}@{}}
\toprule
Metric & Replication (n=20) & Paper (Fig 4 + Table S4) \\
\midrule
Length range & 18,753--164,297 bp & Same order-of-magnitude \\
GC range & 30.66--43.84\% & Consistent \\
MASH pairwise mean $d$ & 0.924 & Fig 4 VIRIDIC heat map is very sparse --- \checkmark \\
Frac pairs $d \approx 1.0$ & 91.05\% & Consistent \\
Frac pairs $d < 0.10$ & 3.68\% & Small ``islands'' in Fig 4 --- \checkmark \\
Total predicted proteins (Prodigal) & 2{,}497 & --- \\
Total protein clusters (MMseqs2 30/50) & 1{,}495 & --- \\
Singleton clusters & 1{,}035 (69.2\%) & --- \\
Clusters shared by $\geq$2 phages & 460 (30.8\%) & --- \\
\textbf{Clusters shared by all 20 (core)} & \textbf{0} & Consistent with paper's ``no strict core genome'' finding \\
Largest cluster size & 9 (out of 20) & --- \\
\bottomrule
\end{longtable}

The 20-focal-phage BIONJ tree (Newick file \texttt{report/evidence/mash20\_nj.nwk}) recovers three quantitatively-visible clades:
\begin{itemize}
    \item \emph{Wbetavirus} sub-clade: \texttt{NC\_007814.1} (Fah), \texttt{DQ222851.1} (Cherry), \texttt{KY963371.1} (Carmel\_SA), \texttt{NC\_048628.1} --- all \emph{B.\ anthracis / B.\ cereus} siphoviruses, matching Table S4 lineage.
    \item SPP1-like sub-clade: \texttt{NC\_020478.1} $\equiv$ \texttt{KC330681.1} $\equiv$ \texttt{NC\_020479.1} $\equiv$ \texttt{NC\_041858.1} --- all near-identical \emph{B.\ subtilis} siphoviruses (branch lengths $<$ 0.01).
    \item Mixed Myoviridae scatter across the rest.
\end{itemize}

This is qualitatively the same clade structure the paper visualises around Figs 5-6.

\section{Per-claim what worked / what didn't}

\begin{longtable}{@{}p{1cm}p{7cm}p{7cm}@{}}
\toprule
Claim & What worked & What didn't (or wasn't attempted) \\
\midrule
C1 & MMseqs2 on 35k proteins recovered a 49\% singleton fraction and 0 clusters with $\geq$100 members --- unambiguously supports the ``large unknown-function fragment'' sub-claim. & WebMGA COG functional annotation of the 6{,}875 cluster reps not re-run (web service). The specific sporulation-regulator / biofilm / virulence-factor cargo sub-claim of C1 is not independently verified. \\
C2 & 231-genome length + GC + MASH heat map recovers exactly the paper's diverse-but-sparse pattern. & PHASTER prophage prediction on 178 host genomes not re-run (web service); therefore the ``clear boundary between prophages and lytic phages'' sub-claim is untested. \\
C3 & MASH pairwise $d$ statistics on all 53,130 pairs quantitatively support ``low genome similarity''. & --- \\
C4 & --- & Lysis-module type-I vs type-II HMM scan on the 231-panel not run; ``influence on cross-species infestation'' inference not tested. \\
\bottomrule
\end{longtable}

\section{Verdict + justification}

\textbf{PARTIAL --- REPLICATED for 3 of 4 testable headline claims.}

C2 (diversity), C3 (low genome similarity), and the ``large unknown-function fraction'' sub-claim of C1 are all quantitatively supported by an independent re-run on 231 of 236 real re-fetched NCBI genomes. Data availability is excellent (all supplementary tables released with full accessions), the pipeline is standard, and the replication produces the same qualitative and near-quantitative pattern. The two remaining sub-claims (prophage-vs-lytic boundary, specific COG cargo, cross-species lysis-module role) require web services (PHASTER, WebMGA) or additional HMM building that is out of the free-tool wave scope. See \texttt{failure\_analysis.md} for a full accounting and \texttt{open\_questions.json} for the 5 concrete follow-on studies.

\section{Open questions (5)}

Grounded in what this replication actually surfaced; see \texttt{report/open\_questions.json} for the full JSON with \texttt{basis} and \texttt{next\_steps} for each.

\begin{itemize}
    \item[\textbf{Q1}] How much of the ``low similarity'' headline is an artefact of accession-collection strategy? Our matrix has 1.11\% pairs at $d < 0.05$ despite different accessions --- dereplicate at 95\% ANI and re-check.
    \item[\textbf{Q2}] Does PHASTER's 20-kb-filtered prophage prediction systematically miss ancient integrated prophages, biasing the ``clear prophage$\leftrightarrow$lytic boundary'' claim? VirSorter2 + geNomad rerun on the 10 host genomes would settle it.
    \item[\textbf{Q3}] Are singleton protein clusters (49\% of 6{,}875) enriched in small $<$50 kb phages? If so, the ``unknown function'' story is partly a small-phage annotation bias, not novel biology.
    \item[\textbf{Q4}] What is the true frequency of the two lysis-module types across the full 231-panel (not just the 3 focal phages), and does the distribution correlate quantitatively with host range? HMMER scan against a two-model library would answer.
    \item[\textbf{Q5}] How stable is the 20-focal-phage tree across methods? Compare our MASH-BIONJ tree to a proteome-presence-absence tree via normalized Robinson-Foulds distance.
\end{itemize}

\section{Reproducibility}

All input data are on NCBI (public, no login). All tools are free and installable from bioconda. Total genome data volume: $\sim$22 MB. Total compute: $\sim$2 min on 32-thread uicgpu (plus $\sim$10 min each for the Marker + Nougat PDF extractions running in parallel). All shell scripts, Python analytics, evidence JSONs, and per-genome length/GC/MASH tables are in \texttt{report/evidence/} and \texttt{work/} in this repository.

\end{document}
